\documentclass[10pt,a4paper]{amsart}
\usepackage[margin=19mm]{geometry}
\usepackage{amsmath,amssymb,mathtools}
\usepackage[scr=rsfs,cal=euler]{mathalfa}
\usepackage{xfrac}
\usepackage[unicode,hidelinks,bookmarks=false]{hyperref}
\newcommand{\ie}{i.e.\ }
\newcommand{\cf}{cf.\ }
\newcommand{\resp}{respectively\ }
\newcommand{\Subsection}{\subsection}
\providecommand{\nobreakdash}{\nobreak\hbox{-}}
\setlength{\emergencystretch}{2em}
\multlinegap0pt
\usepackage{bm}
\usepackage{float}
\usepackage{amssymb}
\usepackage{enumerate}
\newtheorem{lemma}{Lemma}
\title{A spectral Lov\'{a}sz--Simonovits theorem}
\author{Yongtao Li, Lihua Feng, Yuejian Peng}
\date{Source: arXiv:2408.01709v2 (CC BY 4.0)}
\begin{document}
\maketitle

\noindent\emph{Source and licence.} A spectral Lov\'{a}sz--Simonovits theorem. Version arXiv:2408.01709v2 (CC BY 4.0), distributed by arXiv under the Creative Commons Attribution 4.0 International licence, http://creativecommons.org/licenses/by/4.0/, as recorded in the arXiv licence metadata for this exact version and shown on the abstract page https://arxiv.org/abs/2408.01709v2. Full licence notice: \texttt{licenses/arxiv-2408.01709-CC-BY-4.0.txt}.\par\smallskip

\noindent\textit{Editorial scope.} Counted unit: the article's lemma eigen-entry (source0.tex lines 1216--1219). The article's complete original proof of it is delivered with it (source0.tex lines 1222--1288, one \texttt{proof} environment of 67 lines). No mathematics is written, repaired or re-derived: the statement and the proof are the article's own lines, in the article's own notation, and no retained line is substituted or re-flowed.  Retained uncounted as the source-local context the proof needs: lines 1149--1151; lines 1174--1183.  Nothing was omitted from any retained span, so the package carries no omission record.  No retained line contains a citation, so the wrapper carries no bibliography and no bibliography entry.  No retained line contains a figure environment, an image-inclusion command or drawing code, so no image is delivered and the package declares no asset dependency.  Editorial formatting only, in addition: an AMS \texttt{amsart} preamble that restates the article's own declarations and macros used by the retained lines, namely the environments \texttt{lemma} on separate counters, together with the standard packages those lines need.  The retained environments print in this document as Lemma 1 in order of appearance, whereas the article prints the same items with the numbering of its own full document.  The complete native source of the article as distributed in the arXiv e-print for this exact version is bundled unchanged as \texttt{sources/163299/source0.tex}.
\begin{lemma} \label{lem-ST-q}
We have $e(S) + e(T)\le q$. 
\end{lemma}

\begin{lemma} \label{refine-degree}
With the previous notation, we have  
\[ e(S,T) > \frac{n^2}{4}  - 2q   \]
and 
 \[ \frac{n}{2} -  \sqrt{2q} < |S|, |T| 
 < \frac{n}{2} +\sqrt{2q}.  \]
Furthermore, it follows that  
\[  \frac{n}{2} -2q \le \delta(G) \le 
\lambda (G)\le \Delta (G) \le \frac{n}{2} + q + \sqrt{2q}.  \]
\end{lemma}

\begin{lemma} \label{eigen-entry}
For every $u\in V(G)$, we have 
\[  {x}_u > 1- \frac{20q}{n}. \]  
\end{lemma}

\begin{proof} 
Let $z\in V(G)$ be a vertex 
such that ${x}_z$ attains the maximum 
eigenvector entry of $\bm{x}$ and ${x}_z=1$. 
Without loss of generality, we may assume that $z\in S$.  
Our proof of this lemma falls naturally into two steps.  
\medskip 

{\bf Step 1.}   For every $u \in S$, 
we have ${x}_u > 1- \frac{11q}{n}$. 

\medskip 

 Using Lemma \ref{lem-ST-q} gives $e(S)\le q$ and  $d_S(u)\le q $. By Lemma \ref{refine-degree},  we have  
$d_T(u) = d(u)-d_S(u)  \ge ( \frac{n}{2}-2q )-q  = \frac{n}{2}-3q$. 
By Lemma \ref{refine-degree} again, we  get 
$|T|\le \frac{n}{2} + \sqrt{2q}$. Then 
\begin{eqnarray*}
 |N_T(z)| - |N_T(u)\cap N_T(z)|  
 &=& -|N_T(u)| + |N_T(u) \cup N_T(z)|  \\
& \le& - \left(\frac{n}{2} -3q \right) +  
\left(\frac{n}{2}  + \sqrt{2q} \right) \\
&=& 3q + \sqrt{2q} .
\end{eqnarray*}  
Hence, we have 
\begin{eqnarray*}
\lambda {x}_u-\lambda {x}_z 
= \sum_{v\sim u} {x}_v  - \sum_{v\sim z} {x}_v 
\ge  -\sum_{v\in T, v\sim z, v\not\sim u} {x}_v 
-\sum_{v\in S, v\sim z} {x}_v. 
\end{eqnarray*}
Applying the fact ${x}_v \le 1$, it follows that 
\begin{eqnarray*}
\lambda {x}_u-\lambda {x}_z  
\ge  -\Bigl( |N_T(z)| -|N_T(u)\cap N_T(z)| \Bigr) 
- |N_S(z)| 
\ge  - 4q -\sqrt{2q}. 
\end{eqnarray*}
Recall that ${x}_z=1$ 
and $\lambda > {n}/{2}$. Then for each $u\in S$, 
we have 
\begin{equation*}\label{verct1}
{x}_u\ge 1-\frac{4q + \sqrt{2q} }{\lambda }> 1-\frac{8q + 2\sqrt{2q} }{n} > 1-\frac{11q}{n}.
\end{equation*}

{\bf Step 2.}   
For every $u\in T$, we have ${x}_u> 1- \frac{20q}{n}$. 

\medskip 

By the previous step, we get 
 $$\lambda  {x}_u=\sum_{v\in N(u)} {x}_v\ge \sum_{v\in  N_S(u)} {x}_v 
 > \left(1-\frac{8q + 2\sqrt{2q} }{n} \right)d_S(u).$$ 
By  Lemma \ref{refine-degree}, we can see that 
$d(u)\ge \delta (G)\ge \frac{n}{2} - 2q$, 
which together with $d_T(u)\le q$ gives 
$ d_S(u)=d(u)- d_T(u)\ge  \frac{n}{2}-3q$.
 Hence, for $q\ge 2$, we have 
 \begin{eqnarray*}
 {x}_u >  \frac{(1-\frac{8q + 2\sqrt{2q} }{n})d_S(u)}{\lambda } 
 \ge  \frac{(1-\frac{8q + 2\sqrt{2q} }{n})(\frac{n}{2}-3q)}{
 \frac{n}{2}+q + \sqrt{2q}} 
 \ge 1- \frac{20q}{n}.
 \end{eqnarray*} 
Combining the above two claims, 
 the desired result follows immediately.
\end{proof}

\end{document}
